\section{Evaluation: From Task Success to Representation Evidence}
\label{sec:evaluation}

\subsection{Why task success is insufficient}

Closed-loop success is necessary for a control method but weak evidence about a
specific latent mechanism. A policy can infer phase from vision, memorize object
trajectories, exploit a scripted receiver, or rely on privileged object state.
It may succeed without using tactile inputs at all. Conversely, a representation
can predict contact transitions accurately while a poorly tuned controller masks
the benefit. Evaluation should therefore connect four levels: observation
integrity, representation probes, open-loop dynamics, and closed-loop behaviour.

\begin{figure}[t]
\centering
\begin{tikzpicture}[
  rung/.style={draw=reviewblue, fill=reviewlight, minimum height=9mm,
               align=center, font=\small},
  strong/.style={draw=revieworange, fill=reviewwarm, minimum height=9mm,
                 align=center, font=\small},
  arr/.style={-{Latex[length=2mm]}, thick, draw=reviewgray}
]
\node[rung, minimum width=30mm] (integrity) at (0,0) {data integrity\\and leakage checks};
\node[rung, minimum width=30mm] (probe) at (3.4,0.75) {transition, edge,\\load, wrench, slip};
\node[rung, minimum width=30mm] (pred) at (6.8,1.5) {future prediction\\and calibration};
\node[rung, minimum width=31mm] (control) at (10.2,2.25) {closed-loop success,\\safety, recovery};
\node[strong, minimum width=32mm] (causal) at (13.65,3.0) {OOD topology and\\causal interventions};
\draw[arr] (integrity)--(probe);
\draw[arr] (probe)--(pred);
\draw[arr] (pred)--(control);
\draw[arr] (control)--(causal);
\node[font=\scriptsize, text=reviewgray] at (6.8,-0.75) {Increasing strength of evidence for a structured representation claim};
\end{tikzpicture}
\caption{An evidence ladder for representation research. Higher rungs depend on,
rather than replace, the lower-level checks.}
\label{fig:evidence_ladder}
\end{figure}

\subsection{Matched-budget baseline matrix}

The baseline matrix in \cref{tab:baseline_matrix} isolates the source of an
improvement. All learned variants should share the same local tactile encoder,
observations, history, training examples, optimizer budget, and approximately
matched parameter/compute count. If a graph variant receives link poses while a
Transformer does not, the study tests information access, not factorization.
Two additional controls are necessary after the core-paper audit. A predictive
contact-grounding baseline tests whether future tactile/state objectives alone
explain the gain~\cite{heng2025vitacformer,xu2026contactgrounded}; a dynamic taxel
graph tests whether local sensor-level message passing already suffices
~\cite{ning2026touchwgnn}.
The surveillance update adds three more controls: a directed contact--deformation--
slip predictor~\cite{xue2026hitacwam}, training-only temporal-flow grounding
~\cite{yang2026temporalflowvla}, and prediction-error recurrent inference
~\cite{sawada2026predvla}. Cross-embodiment visual prediction such as CLAP is a
useful system-level baseline where video pretraining is available
~\cite{liu2026clap}, but it must receive the same deployment observations when
used to isolate the value of touch.

\begin{table}[t]
\centering
\caption{Minimum matched-budget experiment matrix. The oracle graph is an upper
bound and must never be presented as a deployable method.}
\label{tab:baseline_matrix}
\begin{tabularx}{\linewidth}{P{0.20\linewidth} C C C Y}
\toprule
\textbf{Model} & \textbf{Same local encoder} & \textbf{Temporal context} & \textbf{Explicit structure} & \textbf{Question isolated} \\
\midrule
Proprioception-only & n/a & matched & none & Is touch needed? \\
Tactile concat & \yes & matched & none & Does local touch help? \\
Cross-attention & \yes & matched & implicit & Is flexible token interaction sufficient? \\
Temporal Transformer & \yes & matched & implicit & Is time/context the main source of gain? \\
Prediction-error recurrent & \yes & matched & implicit/predictive & Does observation-driven latent correction explain the gain? \\
Predictive grounding & \yes & matched & implicit/predictive & Do future tactile/state objectives explain the gain? \\
Directed tactile hierarchy & \yes & matched & ordered physical heads & Does contact $\rightarrow$ deformation $\rightarrow$ slip ordering beat flat heads? \\
Training-only temporal grounding & \yes & matched & privileged supervision only & Does richer training supervision, rather than deployed structure, explain the gain? \\
Spatially anchored tokens & \yes & matched & coordinates & Does a shared frame suffice? \\
Dynamic taxel graph & \yes & matched & local learned edges & Does sensor-level graph structure suffice? \\
Static graph & \yes & matched & fixed edges & Does relational routing help without edge inference? \\
Inferred dynamic graph & \yes & matched & observed/inferred & Does topology inference add value? \\
Oracle graph & \yes & matched & privileged edges & What is the attainable structural upper bound? \\
Random / shuffled graph & \yes & matched & invalid control & Does edge semantics, rather than regularization, matter? \\
\bottomrule
\end{tabularx}
\end{table}

\FloatBarrier

For the tool-use phase, contact-centric alternatives are also mandatory: compare
system-level interface edges with a continuous contact-field condition and a
contact-anchored policy frame~\cite{ma2026semanticcontact,cui2026contactanchored}.
These are phase-specific baselines rather than reasons to give the graph model
extra geometry or labels.

A useful diagnostic is the fraction of oracle improvement retained by inferred
structure. Let $M$ be a higher-is-better metric and $M_U$, $M_I$, and $M_O$ the
unstructured, inferred, and oracle scores. Then

\begin{equation}
R_{\mathrm{oracle}}=
\frac{M_I-M_U}{M_O-M_U+\epsilon}.
\label{eq:oracle_retention}
\end{equation}

This ratio is descriptive, not a universal target. It is unstable when the oracle
gain is small and should always be reported with raw scores and uncertainty. For
error metrics, the signs must be adjusted. Its value is that it distinguishes a
weak edge estimator from an unhelpful structural hypothesis.

\subsection{Metrics aligned to the claimed mechanism}

\paragraph{Contact and phase.}
Use per-class precision, recall, and macro-F1 for giver-only, dual-contact,
receiver-only, no-contact, tool-contact, and slip states. Frame accuracy should be
supplemented with transition timing error, segment-level F1, and early-warning
horizon. Long stable segments otherwise dominate the score.

\paragraph{Edges.}
Dynamic-edge evaluation requires precision--recall AUC under class imbalance,
onset/offset timing, false persistent edges, and calibration. Node/edge identity
must be permutation-consistent where appropriate. Testing only edges used to
construct training labels risks circularity; held-out contact configurations and
sensor ablations are stronger.

\paragraph{Load and wrench.}
Report load-share MAE, absolute force/wrench MAE in a stated frame, angular error
for force direction, normalized error across object masses, and horizon-specific
future prediction. Estimated tactile wrench should not be described as independent
ground truth. Gravity compensation, filtering, and coordinate transformations
belong in the method, not supplementary guesswork.

\paragraph{Slip and uncertainty.}
Slip is an event-detection problem with severe imbalance. Precision--recall,
false alarms per minute, and warning time are more informative than accuracy.
Fine-grained contact-state and slip-direction sensing provides a concrete local
classification reference~\cite{ma2026slip}, but local class accuracy should not
be reported as evidence of cross-interface load inference.
For uncertainty, use reliability diagrams, expected calibration error, Brier
score, negative log likelihood, and risk--coverage curves. Calibration should be
re-evaluated under object, sensor, and topology shift.

\paragraph{Duration, relaxation, and damage.}
Sustained-contact evaluation should stratify error by contact duration,
temperature, material, and load/unload history; relaxation-aware force sensing
shows why a single aggregate MAE can hide systematic drift
~\cite{wang2026relaxation}. For deformable objects, report task success and a
separate deformation or damage constraint. SoftVTBench's deformation-aware
success rate is one example of requiring both completion and bounded peak
deformation~\cite{jing2026softvtbench}; the exact threshold and evaluator-only
state must remain explicit.

\paragraph{Closed-loop outcomes.}
Success should be paired with time, peak force, impulse, drop rate, object damage
proxy, path error, number of recovery events, and unnecessary grip. A structured
model that achieves the same success with lower peak force or earlier safe release
can still be valuable.
Failure-aware execution further requires separate accounting for detection,
local retry, reset/recovery, and irrecoverable failure. FLARE operationalizes
these branches with an online monitor and distinct retry/reset mechanisms
~\cite{zhao2026flare}; it is a visual--language recovery framework rather than a
tactile representation, but it makes recovery-aware metrics a necessary control
for long-horizon claims.

\subsection{Splits and out-of-distribution tests}

Random frame splits are invalid for temporally correlated episodes. The split
hierarchy should include, where relevant:

\begin{itemize}
  \item episode and contiguous contact-sequence holds-out;
  \item object identity, geometry, mass, friction, and compliance holds-out;
  \item grasp pose and hand-role swap;
  \item collection session, operator/demonstrator, and calibration session;
  \item sensor instance, sensor type, missing channels, and added latency;
  \item embodiment or morphology;
  \item task/contact topology, such as handover training and tool-use testing.
\end{itemize}

The split must be made before window extraction so adjacent windows cannot cross
train/test boundaries. Hyperparameter selection should use a validation split
with the same unit as the intended generalization claim.

The 30-paper full-text core contained zero studies that met a strict
contact-topology OOD criterion under deployment-observable inference. Strict
topology OOD holds out the contact-participant or edge structure itself (for
example, hand--object--hand training followed by hand--tool--workpiece--hand
testing). Unseen object instances, new tool geometry, contact-mode variation
within a fixed participant structure, and semantic tool--action composition are
valuable tests, but they must be reported separately rather than relabelled as
topology transfer.

\subsection{Causal and counterfactual interventions}

Correlational probes show what can be decoded; interventions show whether the
claimed structure matters. Low-cost interventions include time-shuffling one
hand's tactile stream, swapping left/right roles, masking one fingertip, delaying
receiver contact, perturbing the workpiece, rewiring predicted edges, or freezing
edge probabilities. Physical interventions include known mass/friction changes,
brief stabilizer motion, tool stiffness changes, or an unexpected contact. The
expected direction of change should be stated before the experiment.

For temporal baselines, mask history content, shuffle its order, and disable
observation-driven latent correction exactly rather than replacing the model with
an unrelated architecture. TemporalFlow-VLA and PredVLA make these interventions
especially concrete~\cite{yang2026temporalflowvla,sawada2026predvla}. For
hierarchical tactile prediction, compare flat heads, stop-gradient removal, and
dependency-order shuffling~\cite{xue2026hitacwam}.

For example, delaying receiver contact should delay load-share crossover and safe
release. If a model's phase prediction is unchanged, it may be following a motion
schedule rather than tactile evidence. Rewiring a tool--workpiece edge should
damage future-wrench prediction more than rewiring an inactive relation. If all
rewirings have equal effect, the graph may act as generic regularization.

\subsection{Statistics and reproducibility}

Formal results should use at least five independently seeded runs where training
variance is material; early internal decisions can use three seeds but must not be
reported as definitive. Pair results by the same data split and environment seed.
Report effect sizes and confidence intervals, not only significance tests. For
episode success, use bootstrap or binomial intervals with episode as the unit;
for dense time series, aggregate within episode before testing or use hierarchical
resampling.

Reproducibility requires a data lineage from raw sensor packet to metric: hashes or
version identifiers, synchronization configuration, calibration, split manifests,
normalization fitted only on training data, model configuration, seeds, software
environment, checkpoint selection, and evaluation script. Code audits should
distinguish a paper's claims from what a public snapshot executes, as illustrated
by the PhysGraph case in \cref{sec:structured}.

\begin{table}[H]
\centering
\caption{Claim-to-evidence checklist. A paper need not include every row, but the
evidence should match its stated contribution.}
\label{tab:claim_evidence}
\begin{tabularx}{\linewidth}{P{0.25\linewidth} Y Y}
\toprule
\textbf{Claim} & \textbf{Minimum comparison} & \textbf{Stronger evidence} \\
\midrule
Touch is necessary & proprio/vision-only and tactile ablation & sensor masking and latency interventions \\
Explicit structure helps & matched attention/temporal/static graph & random-edge control and oracle--inferred analysis \\
Edges are observation-grounded & no test-time privileged inputs & pose/sensor corruption and edge calibration \\
Representation is physical & interface/load predictive targets & force-balance residuals and counterfactual response \\
Generalizes & object/session/contact-sequence split & sensor, embodiment, role, and topology holds-out \\
Improves control & multi-seed success and safety metrics & disturbance recovery and risk--coverage analysis \\
\bottomrule
\end{tabularx}
\end{table}

\FloatBarrier
